% 3.2 Recomendación y trabajo a futuro

\section{Recomendación y trabajo a futuro}

En la medida que se ha concluido esta fase final del proyecto, es menester reflexionar sobre los logros alcanzados, así como también identificar aquellas áreas que presentan oportunidades para el crecimiento y mejora continua. En esta sección, se explorarán las áreas identificadas como puntos críticos para el desarrollo continuo, resaltando posibles mejoras, capacidad de innovación y recomendaciones para la realización de un proyecto igual o similar.

Es posible que el área en la que se coloque el dispositivo se trate de un terreno inestable por lo que, para solucionar este problema, se propone que a futuro se desarrolle un sistema que permita nivelar el dispositivo ya sea de forma manual o automatizada.
 
Por otro lado, resultaría bastante satisfactorio poder adaptar el dispositivo para trabajar en intemperies, principalmente en presencia de lluvia.

Una mejora al triturador podría ser útil también, puesto que este presenta ciertas dificultades cuando los vegetales se introducen en tiras largas.

Continuando con las áreas de mejora, resulta prudente mencionar el cómo transportar el dispositivo en un terreno inestable o simplemente trasladarlo una distancia considerable, ya que el compostador es bastante pesado como para estarlo cargando. 

La experiencia del usuario siempre ha sido importante en la realización de este proyecto y aunque se ha buscado cumplir con este propósito, siempre es posible mejorar. Por ejemplo, cambiar el tipo de pantalla a una táctil podría resultar bastante atractivo, refiriéndonos a la parte del monitoreo del proceso, lograr que se pueda acceder de manera totalmente remota, son detalles que a pesar de ser menores, no dejan de ser importantes, y pueden hacer la diferencia en la elección del producto. 

Dos aspectos en los que se considera que el producto podría innovar, sería el manejo de los gases naturales para su uso como energía eléctrica y térmica. El segundo aspecto gira entorno a la descarga del compost, lograr una descarga automatizada podría ser un diferenciador en el mercado.

De la misma forma, la implementación de un sistema de drenaje para la evacuación de todos los lixiviados, producto del proceso de compostaje, permitiría eliminar gran parte de la humedad presente en el compost, lo que en el presente impide obtener un resultado final con un menor porcentaje de agua, obteniendo así una mezcla más parecida a un lodo que a una tierra negra.

Otro punto importante que ayudaría con la disminución del porcentaje de agua presente en el producto final seria el adaptar un mejor sistema de calefacción que permita la eliminación de todos estos lixiviados, o en su defecto considerar una etapa de complementaría externa al sistema de compostaje para la eliminación de estos líquidos.

Pensar en una segunda versión del proyecto que resuelva este problema incluyendo los antes mencionados, podría resultar un tanto ambicioso, sin embargo, se tiene la convicción de que los ingenieros que se forman en esta institución han sido capaces de sorprendernos en más de una ocasión. 

Por otro lado, después de toda la experiencia que ha brindado este proyecto, resulta grato poder compartir algunas recomendaciones para aquellas personas que deseen construir un proyecto similar. Antes de adentrarse en el diseño del dispositivo definir primordialmente como ingresar y extraer la materia. Prestar especial atención al “hermetismo” que se busca pues las fugas contribuyen a la pérdida de calor.

Hacer énfasis en la fuente de calor del compostador, adaptar el mecanismo para que el calor sea transmitido al interior del contenedor puede ser un reto bastante interesante.

\newpage
